\documentclass[11pt]{article}
\usepackage{fontspec}
\usepackage[french]{babel}
\usepackage{amsmath}
\usepackage{amssymb}
\usepackage{amsthm}
\usepackage{mathrsfs}
\usepackage{tikz-cd}
\usetikzlibrary{babel}
\usepackage{enumitem}

\newcommand{\printedpage}[1]{\marginpar{\footnotesize p.~#1}}
\newcommand{\msp}[1]{\marginpar{\footnotesize [#1]}}
\newcommand{\Ch}{\underline{\mathrm{Ch}}}

\begin{document}

\printedpage{4}
% head: A. Grothendieck : Dérivateurs, Ch. XII, Édition des Universités de Montpellier II et Paris VII
On utilise la démonstration de I, p. 75, à l'aide de la catégorie nommée $\Phi(X)$ ($\mathrm{Ob}\,\Phi(X) = \mathrm{Fl}(X)$, mais avec une notion de flèche dans $\Phi(X)$ un peu particulière). On trouve que $f \in W \iff \Phi(f) \in W$, et de même $f^\circ \in W \iff \Phi(f) \in W$.

NB. Dans la démonstration de la proposition 2, comme dans celle de la proposition 1 (pour l'implication b) $\Longrightarrow$ a)), on n'a utilisé la propriété de localisation W(2) que sous la forme faible du corollaire, i.e. sous la forme $f$ $W$-asphérique $\Longrightarrow f \in W$.

\textbf{Corollaire} (de la prop. 2) : \textit{Soit $W \subseteq \mathrm{Fl}(\mathrm{Cat})$ satisfaisant l'axiome} W(1) (saturation). \textit{Alors la conjonction des axiomes} W(2), W(3) \textit{équivaut à celle des axiomes duaux} W($2'$), W($3'$) (\footnote[5]{Itou si on remplace dans cet énoncé W(2) par W(2 bis).}).

Il suffit de prouver que 1,2,3 $\Longrightarrow$ $2'$,$3'$.

\msp{page 4}

Pour $2'$ \texttt{[plutôt} $3'$\texttt{]}, c'est le corollaire 2 ci-dessus. Pour $3'$ \texttt{[plutôt} $2'$\texttt{]}, cela résulte aussitôt de la proposition 2. (NB Cet énoncé resterait valable, si on aurait pris W(2) sous la forme faible seulement, cf. le NB page précédente.)

Je dois encore mettre un axiome de la forme L5, L6 (cf. VI, p. 43-44). Mais l'énoncé de ces axiomes utilise la notion de Hot-fibrés, ou de foncteurs parfaits, i.e. fait appel à la connaissance de $W_\infty$, ce que j'aimerais éviter (si je peux). Je dois donc introduire la notion de $W$-fibration, sur le modèle de ce que j'ai fait dans le contexte des dérivateurs, dans I (cf. p. 111-126). On introduit la notion de \textbf{D}-fibration à la page 117 (prop. 5), par \texttt{[des]} conditions équivalentes, dont certaines ont un sens en termes de la donnée de $W$ sans plus. Soit donc
\[ f : X \longrightarrow Y \]
dans Cat et considérons les conditions

\begin{enumerate}
\item[a)] Si $Y'' \longrightarrow Y'$ dans $\mathrm{Cat}/Y$ est dans $W$, alors $X'' \longrightarrow X'$ aussi (le changement de base par $X \longrightarrow Y$ transforme $W$ dans $W$).
\item[a$'$)] Comme a), mais avec $Y''$, $Y'$ $W$-asphériques.
\item[a$_1$)] Comme a), dans le cas particulier où $Y'$ a un objet final et $Y''$ est la catégorie ponctuelle.
\item[a$_2$)] Dual de a$_1$), avec objet initial de $Y'$.
\end{enumerate}

On a donc
\[
\begin{tikzcd}[row sep=large]
 & & \mathrm{a} \arrow[d, Rightarrow] & \text{$W$-fibration ultraforte} & \\
 & & \mathrm{a}' \arrow[dl, Rightarrow] \arrow[dr, Rightarrow] & \text{$W$-fibration forte} & \\
\text{$W$-fibration gauche} & \mathrm{a}_1 & & \mathrm{a}_2 & \text{$W$-fibration droite}
\end{tikzcd}
\]
% foot: Version 7 décembre 2002
\printedpage{5}
% head: A. Grothendieck : Dérivateurs, Ch. XII, Édition des Universités de Montpellier II et Paris VII
et dans le cas $W = W_\infty$ toutes ces conditions sont équivalentes. J'ai envie \texttt{[de]} prendre la plus forte de ces conditions, mais je dois me donner des moyens pour trouver des $W$-fibrations en un sens aussi fort. Par exemple j'aurais envie de savoir que si $X$
\msp{page 5}
est dans Cat, alors $\Ch_\infty(X) \longrightarrow X \times X$ ou $\Ch(X) \longrightarrow X \times X$ est une $W$-fibration. Dans les deux cas, et (disons) pour $W_\infty$, je le prouve en prouvant que $p$ est à la fois lisse, et propre. Il faudrait donc revoir les notions de lissité et de propreté dans le contexte de la donnée d'un $W$.

Si je me rappelle bien, la propriété a$'$) \texttt{[plutôt a}$_1$\texttt{)]} était quelque chose d'assez aisément vérifiable, il suffisait de supposer que les
\[ X_y \longrightarrow X/y \qquad \text{et} \qquad X/y \longrightarrow X/y' \]
(pour $y \in \mathrm{Ob}\,Y$ resp. $u : y \longrightarrow y'$ dans $\mathrm{Fl}\,Y$) sont dans $W$, et que ça reste vrai par tout changement de base - ce qui revient à la condition a$_1$) sur $f$. On pourrait donc formuler une condition sur $W$ ainsi :

\begin{enumerate}[leftmargin=*]
\item[\textbf{W(4)}] (Axiome des $W$-fibrations (\footnote[6]{Voir variante affaiblie suffisante, dans le NB p. 10. Voir encore [?] de variantes pages 13, 14 ci-dessous.}).) Soit $f : X \longrightarrow Y$ dans Cat satisfaisant la condition de $W$-fibration a$_1$) ci-dessus (p. 4), alors elle satisfait même à a), i.e. le changement de base par $f$ transforme flèches dans $W$ en flèches dans $W$ (\footnote[7]{i.e. une $W$-fibration à gauche est une fibration ultraforte.}).
\end{enumerate}

Grâce à la proposition 2, il s'ensuit que la condition duale a$_2$) sur $f$ implique la condition de $W$-fibration forte, donc l'axiome W(4), en apparence asymétrique, signifie aussi que les quatre conditions de fibration sur $f$ sont équivalentes, il est donc en fait symétrique.

\textbf{NB} Je ne vois pas, sans axiomes supplémentaires sur $W$, comment passer de a$_1$ à a$'$ seulement. La difficulté, c'est de montrer, moyennant a$_1$ seulement, que si $Y$ est $W$-asphérique, l'inclusion $X_y \longrightarrow X$ (où $y \in \mathrm{Ob}\,Y$) est dans $W$. Dans le contexte des dérivateurs \textbf{D}, cela se voyait par voie \textbf{D}-cohomologique.

Cet axiome très fort suffit à impliquer que si $X$, $Y$ sont $W$-fibrés sur $S$, et si $f : X \longrightarrow Y$ est un $S$-morphisme qui est dans $W$, alors $f_s : X_s \longrightarrow Y_s$ est dans $W$ ($\forall\ s \in S$). Montrons déjà au moins l'inverse.

\msp{page 6}

\textbf{Proposition 3.} \textit{Soient $X$, $Y$ dans} Cat$/S$\textit{, $f : X \longrightarrow Y$ une $S$-flèche. Supposons ($X$, $Y$ $W$-fibrés) que $\forall\ s \in S$, $f_s : X_s \longrightarrow Y_s$ soit dans $W$, alors $f \in W$} (\footnote[8]{\textit{Sans} axiome W(4). C'est la première fois qu'on a besoin de W(2) sous sa forme forte, à l'exclusion de W(2 bis).}).

Il s'impose d'utiliser l'axiome de localisation, donc de regarder les $f/s$ induits par $f$
% foot: Version 7 décembre 2002
\printedpage{6}
% head: A. Grothendieck : Dérivateurs, Ch. XII, Édition des Universités de Montpellier II et Paris VII
\[
\begin{tikzcd}[row sep=large]
X/s \arrow[r, "f/s"'] & Y/s \\
X_s \arrow[u, "i_s^X"] \arrow[r, "f_s"] & Y_s \arrow[u, "i_s^Y"'] \quad .
\end{tikzcd}
\]

La condition a$_1$ suffit à impliquer que $i_s^X$ et $i_s^Y$ sont dans $W$, donc $f/s$ est dans $W$ si et seulement si $f_s$ l'est. On n'a pas eu à utiliser l'axiome W(4), et il suffisait de supposer que $X$, $Y$ soient fibrés ‘$W$-cohomologiques’ (ou $W$-fibrés à gauche) sur $S$, ou dualement, des $W$-fibrés à droite.

De plus :

\textbf{Corollaire} (\footnote[9]{Sans axiome W(4).}). \textit{Soit $S_0 \subseteq \mathrm{Ob}\,S$ une partie de $S$ telle que $S_0 \longrightarrow \pi_0(S)$ surjectif. Alors dans la proposition 3 il suffit de supposer $f_s \in W$ pour $s \in S_0$.}

Il suffit de voir que si $u : s \longrightarrow s'$ est une flèche dans $S$, alors $f_s \in W \iff f_{s'} \in W$. On utilise le fait que $f_s \in W \iff f/s \in W$, et de même pour $f_{s'}$. Donc à voir que $f/s \in W \iff f/s' \in W$. Or on a
\[
\begin{tikzcd}[row sep=large]
X/s \arrow[r, "f/s"] \arrow[d, "X/u"'] & Y/s \arrow[d, "Y/u"] \\
X/s' \arrow[r, "f/s'"] & Y/s' \quad ,
\end{tikzcd}
\]
où $X/u$ et $Y/u$ sont dans $W$, d'où le résultat.

Mais revenons à la réciproque : si $f \in W$, alors a-t-on $f_s \in W_S$? \texttt{[plutôt} $f_s \in W$\texttt{?]} Il y a un argument dans ce sens dans VI, p. 23, que je vais reprendre. L'idée est celle-ci. On considère la composante connexe $Y_0$ de $Y$ \texttt{[plutôt} $S_0$ \texttt{de} $S$\texttt{]} contenant $s$, et on construit un
\msp{page 7}
$W$-fibré (au \textit{sens fort} a))
\[ S' \longrightarrow S_0 \subseteq S \]
avec $S'$ \textit{$W$-asphérique.} Comme $X \longrightarrow Y$ sur $S$ est dans $W$, il en sera donc de même de $X' \longrightarrow Y'$ sur $S'$. Comme les fibres de $X'$, $Y'$ en $s'$ s'identifient à celles de $X$, $Y$ en l'image de $s'$ dans $S$, on est ramené au cas où $S$ est $W$-asphérique. Mais alors on a
\[
\begin{tikzcd}[row sep=large]
X \arrow[r, "f"] & Y \\
X_s \arrow[u, "i_s^X"] \arrow[r, "f_s"] & Y_s \arrow[u, "i_s^Y"'] \quad ,
\end{tikzcd}
\]
% foot: Version 7 décembre 2002
\printedpage{7}
% head: A. Grothendieck : Dérivateurs, Ch. XII, Édition des Universités de Montpellier II et Paris VII
où $i_s^X$, $i_s^Y$ sont dans $W$, car déduites de $\{s\} \longrightarrow S$ (qui est dans $W$, $S$ étant $W$-asphérique) par les changements de base $W$-fibrants $X \longrightarrow S$, $Y \longrightarrow S$. Donc $f \in W \implies f_s \in W$.

NB On n'a utilisé la propriété de fibration pour $X$, $Y$ qu'au sens a$'$) (\footnote[10]{Et même, dans a$'$) on pourrait exiger que $Y'' = e$, $Y'$ $W_\omega$-asphérique, où $W_\omega$ \texttt{[est]} défini page 14, $W_\omega \subseteq W \cap W_\infty \ldots$}).

Il faut donc prouver, si on peut, le

\textbf{Lemme} (?) (\footnote[11]{Est prouvé moyennant l'axiome W(4) (+ des petites vérifications ultérieures).}). \textit{Soit $S$ un objet} $0$\textit{-connexe de} Cat\textit{. Alors il existe un objet $W$-asphérique $S'$ de} Cat\textit{, et une $W$-fibration (au sens} fort \textit{de a), page 4) $S' \longrightarrow S$.}

L'idée naturelle, c'est de prendre, pour $s_0 \in \mathrm{Ob}\,S$ fixé,
\[ S' = \Ch(S; s_0, -) \qquad \text{ou} \qquad S' = \Ch_\infty(S; s_0, -), \]
et la flèche
\[ \beta : c \longmapsto \text{but } c : S' \longrightarrow S. \]
On sait que $S'$ est $W_\infty$-aphérique et que $\beta$ est $W_\infty$-fibrante (et même parfaite). Il faut voir que $S'$ est aussi $W$-asphérique, et $\beta$ $W$-fibrante. Il faut donc revoir la démonstration dans le cas de $W_\infty$, et prouver qu'elle se transforme à $W$, % [sic] aphérique -> asphérique, lettre s omise
\msp{page 8}
moyennant les axiomes qu'il faut sur $W$. J'ai envie de regarder d'abord le choix $S' = \Ch(S; s_0, -)$, sans passage à une limite inductive qui m'obligerait à mettre un axiome de limites sur $W$. Les propriétés homotopiques délicates et essentielles de $\Ch(S)$ sont établies dans VII, p. 157-169, savoir que dans le diagramme
\[ S \xrightarrow{i_S} \Ch(S) \xrightarrow{p_S} S \times S, \]
$i_S \in W_\infty$ et $p_S$ est une $W_\infty$-fibration (en fait, \textit{propre} et \textit{lisse} pour $W_\infty$) (\footnote[12]{Dans l'énoncé des ‘propriétés essentielles’ de $\Ch(S)$, on peut travailler avec $W_\omega$, la plus petite partie de Fl(Cat) satisfaisant les axiomes W(1), W(2 bis), W(3) (cf. page 14), donc $W_\omega \subseteq W \cap W_\infty$.}). Le fait que $i_S \in W$ \texttt{[plutôt} $i_S \in W_\infty$\texttt{]} implique que les foncteurs source et but
\[ \begin{tikzcd} \Ch(S) \arrow[r, shift left, "\sigma"] \arrow[r, shift right, "\beta"'] & S \end{tikzcd} \qquad \sigma = \mathrm{pr}_1 \circ p_S,\ \beta = \mathrm{pr}_2 \circ p_S \]
sont dans $W_\infty$ (puisque $\sigma i_S = \beta i_S = \mathrm{id}_S$). Comme ce sont des Hot-fibrations, il en résulte que les fibres $\Ch(S; s, -)$ resp. $\Ch(S; -, s)$ sont Hot-asphériques. Mais même en admettant que $\sigma$, $\beta$ sont dans $W$ et sont des $W$-fibrations, je ne peux en conclure tout de suite que les fibres sont $W$-asphériques - il y aurait cercle vicieux! Il me faut donc une démonstration directe.

Considérons le foncteur
\[
\begin{tikzcd}[row sep=huge]
\Ch(S; s, -) \arrow[d, "p"'] \\
\underline{Ty} \quad ,
\end{tikzcd}
\]
% foot: Version 7 décembre 2002
\printedpage{8}
% head: A. Grothendieck : Dérivateurs, Ch. XII, Édition des Universités de Montpellier II et Paris VII
c'est un foncteur Cat-fibrant, donc $W$-lisse (je devrai développer la notion de $W$-lissité (\footnote[13]{$f : X \longrightarrow Y$ est $W$-lisse si les $X_y \longrightarrow y\backslash X$ sont $W$-asphériques.}) avec soin, j'admettrai pour le moment les propriétés formelles familières, et les justifierai par la suite). À cause de la $W$-lissité, pour prouver $p \in W$, il suffit de
\msp{page 9}
prouver que les fibres de $p$ sont $W$-asphériques (\footnote[14]{On utilise ici W(2 bis), pas $W(2)$.}). Or les fibres sont les $Z_\tau = \Ch_\tau(S; s, -)$, qui sont des catégories \textit{contractiles} (i.e. $Z_\tau \longrightarrow e$ est un homotopisme), donc $W$-asphériques par W(3) (cf. cor. 3, page 3).

D'autre part, si on admet que $\Ch(S) \xrightarrow{p_S} S{\times}S$ est $W$-lisse (et $W$-propre), donc $\Ch(S) \xrightarrow{\sigma_S} S$ aussi, étant à fibres $W$-asphériques comme on vient de voir, il s'ensuit que $\sigma_S \in W$, d'où $S \xrightarrow{i_S} \Ch(S)$ aussi dans $W$ (puisque $\sigma_S \circ i_S = \mathrm{id}_S$). Donc il reste à prouver seulement que
\[ \Ch(S) \xrightarrow{p_S} S \times S \tag*{$(*)$} \]
est $W$-propre et $W$-lisse. Il en résultera que
\[ \Ch(S; s_0, -) \xrightarrow{\beta} S \]
est lui aussi $W$-propre et $W$-lisse, car déduit de $p_S$ par le changement de base
\[ S \longrightarrow S \times S \qquad s \longmapsto (s_0, s). \]

Il faut ensuite voir que $W$-lisse et $W$-propre implique la propriété de $W$-fibration \textit{sous sa forme la plus forte} a) (alors que je prévois qu'a priori on ne l'aura que sous les formes faibles a$_1$, a$_2$ - pas même sous la forme a$'$). C'est ici donc que l'axiome W(4) des $W$-fibrations sera essentiel (\footnote[15]{sous la forme affaiblie W(4a$''$) (p. 13).})!

Voici l'argument pour la lissité de $(*)$, dans loc. cit. p. 161 ff. Il faut prouver que pour $c \in \mathrm{Ob}\,\Ch(S)$, et une flèche $u \in \mathrm{Fl}(S \times S)$ \texttt{[plutôt} $(u,v) \in$ \texttt{Fl(}$S \times S$\texttt{)]}, une certaine catégorie de ‘relèvements de $u$ de source $c$ ’ \texttt{[plutôt ‘relèvements de} $(u,v)$ \texttt{de but} $c$ \texttt{’]}, $\mathcal{C} = \mathcal{C}(c,u,v)$, est $W$-asphérique (\footnote[16]{En fait, on prouve que $\mathcal{C}(c,u,v)$ est \textit{contractile}, donc elle est $W_\omega$-asphérique.}). Pour ceci, on la compare à une catégorie $\mathcal{C}'$ à objet final, donc contractile et a fortiori $W$-asphérique, à l'aide de deux foncteurs
\[ \begin{tikzcd} \mathcal{C} \arrow[r, shift left, "\varphi"] & \mathcal{C}', \arrow[l, shift left, "\psi"] \end{tikzcd} \]
et des homotopies $\varphi\psi \sim \mathrm{id}_{\mathcal{C}'}$ (sans mérite, $\mathcal{C}'$ étant contractile)
\msp{page 10}
et $\psi\varphi \sim \mathrm{Id}_{\mathcal{C}}$ (on trouve même une flèche $\psi\varphi \longrightarrow \mathrm{id}_{\mathcal{C}}$). Ainsi $\varphi$ est un homotopisme, donc une $W$-équivalence (cor. 3, p. 3), donc $\mathcal{C}$ est $W$-asphérique, $\mathcal{C}'$ l'étant. La propreté de $p_S$ équivaut à la lissité de $p_{(S^\circ)}$, elle est donc prouvée du même coup.
% foot: Version 7 décembre 2002

\end{document}
